\chapter*{РЕФЕРАТ}
\addcontentsline{toc}{chapter}{РЕФЕРАТ}
\noindent\textbf{Ключевые слова:} ансамбли деревьев решений, случайный лес, градиентный бустинг, сдвиг распределения, обобщение на новые домены, табличные данные.

Объект исследования --- ансамбли деревьев решений для классификации и регрессии табличных данных. Предмет исследования --- методы повышения устойчивости этих ансамблей к изменению распределения между исходными и целевыми доменами.

Цель работы --- экспериментально сопоставить существующие методы при выборе модели только по исходным доменам и установить границы обоснованных выводов об их качестве. Рассматриваются IRF, CRF, интерпретация Meta-forests, MaxRM, Era splitting и Anchor boosting, а также сильные базовые ансамбли.

В отчёте разграничиваются эффект настройки гиперпараметров и вклад исследуемого механизма. Приводятся результаты предварительных серий и нового теста с заранее зафиксированным протоколом. Положительные, отрицательные и неполученные статистические результаты сохраняются; универсальное превосходство методов не предполагается.
